\section{Distillation Training Methodology}
\label{sec:distillation-training-methodology}

Distillation is implemented as a supervised fine-tuning task using cross-entropy loss.

\subsection{Cross-Entropy and Log-Probabilities}

Training focuses on Next-Token Prediction. The model is penalized based on how much probability it assigns to the ``correct'' token (the one generated by the teacher).

\begin{itemize}
    \item \textbf{Log-Probability:} Measures the model's confidence in the correct next token.
    \item \textbf{Cross-Entropy Loss:} The negative average of these token log-probabilities. A value closer to 0 indicates the student is highly confident in the teacher's reasoning path.
\end{itemize}

\subsection{Answer-Only Loss}

A critical technical detail in distillation is the use of Answer-Only Loss.

\begin{itemize}
    \item \textbf{Mechanism:} The loss is computed only on the tokens generated by the teacher, not the tokens in the prompt.
    \item \textbf{Logic:} The model's task is to generate the answer conditioned on the prompt. Penalizing the model for reproducing the prompt tokens (which are already provided as input) is counterproductive.
\end{itemize}

\subsection{The Training Loop}

Distillation training iterates over the dataset for multiple epochs.

\begin{itemize}
    \item \textbf{Epoch:} One complete pass through the training data. Shuffling the data each epoch helps the model generalize.
    \item \textbf{Metrics:} Progress is tracked via Training Loss (measured on optimized samples) and Validation Loss (measured on a held-out set). Validation loss is the more reliable metric for assessing whether the student's improvements will generalize to new problems.
\end{itemize}
